% ============================================================
\chapter{Mod\`eles AR, MA, ARMA}
\label{ch:arma}
% ============================================================

\section{Exemple motivant}
Les log-rendements journaliers du CAC~40 montrent une corr\'elation faible mais significative aux premiers d\'ecalages. Un mod\`ele AR(1) ou ARMA(1,1) capture cette d\'ependance lin\'eaire.

\section{Processus autor\'egressifs AR$(p)$}\index{AR}

\begin{definition}[AR$(p)$]
\[
X_t = \phi_1 X_{t-1} + \phi_2 X_{t-2} + \cdots + \phi_p X_{t-p} + \eps_t,
\]
o\`u $\eps_t\sim\mathrm{BB}(0,\sigma^2)$. En notation op\'erateur : $\phi(B)X_t = \eps_t$ avec $\phi(z)=1-\phi_1 z-\cdots-\phi_p z^p$.
\end{definition}

\begin{theoreme}[Stationnarit\'e d'un AR$(p)$]\label{thm:ar-stat}
L'AR$(p)$ admet une unique solution stationnaire si et seulement si toutes les racines de $\phi(z)=0$ sont \textbf{en dehors du disque unit\'e} : $|z_i|>1$ pour tout $i$.
\end{theoreme}

\begin{proof}[Esquisse pour AR(1)]
$X_t = \phi X_{t-1}+\eps_t$. Par r\'ecurrence : $X_t = \sum_{j=0}^{\infty}\phi^j\eps_{t-j}$. Cette s\'erie converge dans $L^2$ ssi $\sum\phi^{2j}<\infty$, i.e.\ $|\phi|<1$. La racine de $\phi(z)=1-\phi z$ est $z=1/\phi$, qui est hors du disque ssi $|\phi|<1$.
\end{proof}

\begin{theoreme}[ACF d'un AR(1)]
Si $X_t = \phi X_{t-1}+\eps_t$ avec $|\phi|<1$ :
\[
\gamma(h) = \frac{\sigma^2\phi^{|h|}}{1-\phi^2}, \quad \rho(h)=\phi^{|h|}.
\]
L'ACF d\'ecro\^it exponentiellement.
\end{theoreme}

\begin{proof}
$\gamma(h)=\E[X_t X_{t+h}]=\phi\gamma(h-1)$ pour $h\geq 1$ (par multiplication de l'\'equation AR par $X_{t+h}$ et prise d'esp\'erance). Avec $\gamma(0)=\sigma^2/(1-\phi^2)$.
\end{proof}

\begin{theoreme}[PACF d'un AR$(p)$]
Pour un AR$(p)$ : $\alpha(h)=0$ pour tout $h>p$. La PACF tronque \`a l'ordre $p$.
\end{theoreme}

\section{\'Equations de Yule--Walker}\index{Yule--Walker}

\begin{definition}
Pour un AR$(p)$ stationnaire, les relations :
\[
\gamma(h) = \phi_1\gamma(h-1)+\cdots+\phi_p\gamma(h-p) \quad (h\geq 1)
\]
donnent le syst\`eme de \textbf{Yule--Walker} :
\[
\begin{pmatrix}\gamma(0)&\gamma(1)&\cdots&\gamma(p-1)\\\gamma(1)&\gamma(0)&\cdots&\gamma(p-2)\\\vdots&&\ddots&\vdots\\\gamma(p-1)&\gamma(p-2)&\cdots&\gamma(0)\end{pmatrix}
\begin{pmatrix}\phi_1\\\phi_2\\\vdots\\\phi_p\end{pmatrix}
=\begin{pmatrix}\gamma(1)\\\gamma(2)\\\vdots\\\gamma(p)\end{pmatrix}.
\]
\end{definition}

\section{Processus moyenne mobile MA$(q)$}\index{MA}

\begin{definition}[MA$(q)$]
\[
X_t = \eps_t + \theta_1\eps_{t-1} + \cdots + \theta_q\eps_{t-q} = \theta(B)\eps_t,
\]
o\`u $\theta(z)=1+\theta_1 z+\cdots+\theta_q z^q$.
\end{definition}

\begin{theoreme}[Stationnarit\'e d'un MA$(q)$]
Tout MA$(q)$ est \textbf{toujours stationnaire} (somme finie de bruits blancs).
\end{theoreme}

\begin{theoreme}[ACF d'un MA$(q)$]
\[
\gamma(h) = \begin{cases}\sigma^2\sum_{j=0}^{q-|h|}\theta_j\theta_{j+|h|} & \text{si }|h|\leq q,\\0 & \text{si }|h|>q,\end{cases}
\]
avec $\theta_0=1$. L'ACF \textbf{tronque} au d\'ecalage $q$.
\end{theoreme}

\begin{definition}[Inversibilit\'e]\index{inversibilite@inversibilit\'e}
Un MA$(q)$ est \textbf{inversible} si toutes les racines de $\theta(z)=0$ sont en dehors du disque unit\'e. Dans ce cas, $X_t$ admet une repr\'esentation AR$(\infty)$ : $\pi(B)X_t=\eps_t$ avec $\pi(z)=\theta(z)^{-1}$.
\end{definition}

\section{Processus ARMA$(p,q)$}\index{ARMA}

\begin{definition}[ARMA$(p,q)$]
\[
\phi(B)X_t = \theta(B)\eps_t, \quad\text{i.e.}\quad X_t - \sum_{i=1}^p\phi_i X_{t-i} = \eps_t + \sum_{j=1}^q\theta_j\eps_{t-j}.
\]
\end{definition}

\begin{theoreme}[Conditions de stationnarit\'e et inversibilit\'e]
\begin{itemize}
\item Stationnaire $\iff$ racines de $\phi(z)=0$ hors du disque unit\'e.
\item Inversible $\iff$ racines de $\theta(z)=0$ hors du disque unit\'e.
\item On exige toujours que $\phi(z)$ et $\theta(z)$ n'aient \textbf{pas de racine commune} (sinon le mod\`ele est sur-param\'etr\'e).
\end{itemize}
\end{theoreme}

\begin{theoreme}[Repr\'esentation MA$(\infty)$]
Si l'ARMA$(p,q)$ est stationnaire : $X_t = \psi(B)\eps_t$ avec $\psi(z)=\theta(z)/\phi(z)=\sum_{j=0}^{\infty}\psi_j z^j$. Les $\psi_j$ sont les \textbf{poids de la r\'eponse impulsionnelle}.
\end{theoreme}

\section{Signatures ACF/PACF}

\begin{center}
\begin{tabular}{lcc}
\toprule
\textbf{Mod\`ele} & \textbf{ACF} & \textbf{PACF} \\
\midrule
AR$(p)$ & D\'ecroissance exponentielle/sinuso\"idale & Tronque apr\`es $p$ \\
MA$(q)$ & Tronque apr\`es $q$ & D\'ecroissance exponentielle/sinuso\"idale \\
ARMA$(p,q)$ & D\'ecroissance apr\`es $q$ & D\'ecroissance apr\`es $p$ \\
\bottomrule
\end{tabular}
\end{center}

\section{Impl\'ementation}

\begin{codeblock}[title=\textbf{Python}]
\begin{minted}{python}
import numpy as np
import matplotlib.pyplot as plt
from statsmodels.tsa.arima_process import ArmaProcess
from statsmodels.graphics.tsaplots import plot_acf, plot_pacf

# Simuler un AR(2)
np.random.seed(42)
ar_params = np.array([1, -0.75, 0.25])  # 1 - 0.75*B + 0.25*B^2
ma_params = np.array([1])
ar2 = ArmaProcess(ar_params, ma_params)
x = ar2.generate_sample(nsample=500)

fig, axes = plt.subplots(1, 3, figsize=(15, 4))
axes[0].plot(x, lw=0.7)
axes[0].set_title('AR(2): $\\phi_1=0.75, \\phi_2=-0.25$')
plot_acf(x, lags=25, ax=axes[1], title='ACF')
plot_pacf(x, lags=25, ax=axes[2], title='PACF')
plt.tight_layout()
plt.savefig('ch03_arma.pdf')
plt.show()

# Simuler un MA(2)
ma_params = np.array([1, 0.6, -0.3])
ma2 = ArmaProcess(np.array([1]), ma_params)
y = ma2.generate_sample(nsample=500)

# Simuler un ARMA(1,1)
arma11 = ArmaProcess(np.array([1, -0.8]), np.array([1, 0.4]))
z = arma11.generate_sample(nsample=500)

# Yule-Walker
from statsmodels.regression.linear_model import yule_walker
rho, sigma = yule_walker(x, order=2)
print(f"Yule-Walker: phi = {rho}")
\end{minted}
\end{codeblock}

\section{Exercices}

\begin{exercice}[$\star$]
Calculer l'ACF d'un MA(1) : $X_t = \eps_t+0{,}5\eps_{t-1}$. V\'erifier num\'eriquement.
\end{exercice}

\begin{exercice}[$\star\star$]
Montrer que l'AR(1) $X_t = \phi X_{t-1}+\eps_t$ avec $|\phi|<1$ a une repr\'esentation MA$(\infty)$ : $X_t = \sum_{j=0}^{\infty}\phi^j\eps_{t-j}$.
\end{exercice}

\begin{exercice}[$\star\star$]
Montrer qu'un ARMA(1,1) avec $\phi=\theta$ se r\'eduit \`a un bruit blanc (racine commune).
\end{exercice}

\begin{exercice}[$\star\star$]
D\'eriver les \'equations de Yule--Walker pour un AR(2) et r\'esoudre le syst\`eme.
\end{exercice}

\begin{exercice}[$\star\star\star$ -- Projet]
Simuler des processus AR(1), AR(2), MA(1), MA(2), ARMA(1,1) et ARMA(2,1). Pour chacun, tracer l'ACF et la PACF th\'eoriques et empiriques. Construire un guide visuel d'identification.
\end{exercice}

\begin{formules}
\begin{align*}
\text{AR}(p) &: \phi(B)X_t=\eps_t, \quad \text{stationnaire ssi racines de }\phi(z)\text{ hors disque} \\
\text{MA}(q) &: X_t=\theta(B)\eps_t, \quad \text{toujours stationnaire} \\
\text{ACF AR}(1) &: \rho(h)=\phi^{|h|} \\
\text{ACF MA}(q) &: \rho(h)=0\text{ pour }|h|>q
\end{align*}
\end{formules}
